\documentclass[11pt,a4paper]{article}
\usepackage[margin=2.2cm]{geometry}
\usepackage{amsmath,amssymb,amsthm}
\usepackage{bm}
\usepackage{microtype}
\usepackage{enumitem}
\setlength{\parskip}{2pt}
\setlist{nosep}

\title{Methodology: MAPPO with FOFE Observation Encoding}
\author{}
\date{}

\begin{document}
\maketitle

\section{Methodology}

\subsection{Problem Formulation}
The striker--jammer cooperative engagement is modelled as a Dec-POMDP
$\langle \mathcal{N},\mathcal{S},\{\mathcal{A}_i\},\{\mathcal{O}_i\},
P,\mathcal{R},\gamma\rangle$ with two roles $\rho \in \{s,j\}$ and
$\mathcal{N}=\mathcal{N}_s\cup\mathcal{N}_j$, $|\mathcal{N}_s|=n_s$,
$|\mathcal{N}_j|=n_j$. Each agent emits a factored discrete action
\begin{equation}
\bm{a}_i^t = (a_i^{t,1},\dots,a_i^{t,K}) \in \prod_{k=1}^{K} \mathcal{A}^{(k)},
\qquad
\pi_\rho(\bm{a}\mid o) = \prod_{k=1}^{K} \pi_\rho^{(k)}(a^{(k)}\mid o),
\end{equation}
i.e.\ $K$ independent categoricals per step (a factorised \emph{MultiCategorical}).
Training follows the standard CTDE paradigm: decentralised execution from
$o_i^t$, centralised critic evaluation from the global state $s^t$.

\subsection{Multi-Agent PPO}
We optimise the clipped PPO objective~\cite{schulman2017ppo} in the MAPPO
variant~\cite{yu2022mappo}, with parameter sharing inside each role:
$\pi_{\theta_\rho}$ for role $\rho\in\{s,j\}$, and role-specific centralised
critics $V_{\phi_\rho}$.

\paragraph{Advantage estimation.}
With rollouts of length $T$ collected from $B$ parallel environments, per-role
advantages are computed by GAE~\cite{schulman2016gae}:
\begin{equation}
\delta_t^{\rho} = r_t^{\rho} + \gamma\, V_{\phi_\rho}(s_{t+1})(1-d_t) - V_{\phi_\rho}(s_t),
\qquad
\hat{A}_t^{\rho} = \sum_{l=0}^{T-t-1}(\gamma\lambda)^{l}\,\delta_{t+l}^{\rho},
\end{equation}
with returns $\hat{R}_t^\rho=\hat{A}_t^\rho+V_{\phi_\rho}(s_t)$.
We use $\gamma=0.99$, $\lambda=0.95$.

\paragraph{Clipped policy loss.}
For each role, letting
$r_t(\theta)=\exp\!\big(\log\pi_{\theta_\rho}(\bm{a}_t\mid o_t)-\log\pi_{\theta_{\rho,\text{old}}}(\bm{a}_t\mid o_t)\big)$,
\begin{equation}
\mathcal{L}^{\text{PPO}}_\rho(\theta)
= -\,\mathbb{E}_t\!\left[\min\!\big(r_t(\theta)\tilde{A}_t^\rho,\;
    \mathrm{clip}(r_t(\theta),1\!-\!\epsilon,1\!+\!\epsilon)\,\tilde{A}_t^\rho\big)\right]
- \beta\,\mathbb{E}_t[\mathcal{H}(\pi_{\theta_\rho})],
\end{equation}
with clip range $\epsilon=0.2$ and entropy coefficient $\beta=0.01$.
$\tilde{A}_t^\rho$ denotes the per-minibatch standardised advantage,
$\tilde{A}=(\hat{A}-\mu_{\hat A})/\sigma_{\hat A}$.

\paragraph{Value loss.}
\begin{equation}
\mathcal{L}^{V}_\rho(\phi)=\mathbb{E}_t\!\left[\big(V_{\phi_\rho}(s_t)-\hat{R}_t^\rho\big)^{2}\right].
\end{equation}

\paragraph{Reward normalisation.}
A running second moment over per-step team-mean rewards $\bar r_t=\tfrac{1}{|\mathcal N|}\sum_i r_{i,t}$
is tracked with a Welford estimator; rewards are then divided (no mean
subtraction) by $\max(\sigma,\varepsilon)$ with $\varepsilon=10^{-2}$
before advantage computation.

\paragraph{Adaptations from standard MAPPO.}
\begin{enumerate}[leftmargin=1.25em]
\item \emph{Role-factorised critic.} Instead of a single shared centralised
critic, we instantiate one actor $\pi_{\theta_\rho}$ \emph{and} one critic
$V_{\phi_\rho}$ per role, each optimised by a dedicated Adam instance. The
critics share the same global-state input but learn role-specific value
functions, which we found necessary because the two roles receive different
reward mixtures (Sec.~\ref{subsec:rewards}).
\item \emph{Unclipped value loss.} We omit the PPO value clip; the plain MSE
in~(4) was empirically stable in our setup and avoided a known pathology
where clip-range mis-calibration freezes the critic early in training.
\item \emph{Std-only reward scaling.} We divide rewards by the running std
without mean subtraction, preserving the sign/zero of sparse terminal
rewards (kill / mission-complete), which would otherwise drift to a
non-zero baseline under a Welford-mean subtraction.
\item \emph{Factorised MultiCategorical} action distribution with logit
sanitisation $\mathrm{clip}(\ell,-20,20)$ and NaN guarding, to tolerate
occasional extreme logits produced during the early exploration phase.
\item \emph{Per-role gradient clipping.} Actor/critic gradients are clipped
independently at $\lVert g\rVert_2\le 1$ before each optimiser step.
\end{enumerate}

\subsection{Flexible Observation Feature Encoding (FOFE)}
The observation encoder replaces a fixed top-$K$ flat vector with a
permutation-invariant set encoder applied per observation channel, following
the Sequence-Exchangeable Encoding (SEE) block of Wang et al.~\cite{wang2026fofe}.
Each channel $c\in\{\text{ag},\text{tg},\text{rd}\}$ provides a variable-size
entity set $\mathcal{X}^{(c)}\!=\!\{x_1,\dots,x_n\}$ with a visibility mask
$m\in\{0,1\}^n$.

\paragraph{SEE layer.}
For input $X\in\mathbb{R}^{n\times d_{\text{in}}}$ and hidden width $h$,
\begin{align}
\tilde X &= \mathrm{LayerNorm}(X),\\
E &= \tilde X W_e,\qquad Z=\tilde X W_g,\qquad G=\mathrm{SiLU}(Z),\\
E_{\text{agg}} &= \max_{i:\,m_i=1} E_{i,:} \;\in\mathbb R^{h} \quad\text{(masked max-pool),}\\
\mathrm{SEE}(X,m)_{i,:} &= \big[\,G_{i,:}\odot E_{i,:}\,\big\Vert\,G_{i,:}\odot E_{\text{agg}}\,\big]\in\mathbb R^{2h},
\end{align}
with layer-norm omitted in the first SEE block. The left branch carries
local per-entity context, the right branch injects a gated global summary;
masked rows are zeroed after the concat.

\paragraph{FOFE block.}
A channel encoder stacks $L$ SEE layers, applies a shared per-entity MLP
$f_\psi$, then collapses the set by masked max-pool:
\begin{equation}
\mathrm{FOFE}(\mathcal X,m)
= \max_{i:\,m_i=1} f_\psi\!\big(\mathrm{SEE}_L\circ\cdots\circ\mathrm{SEE}_1(X,m)\big)_i
\;\in\mathbb R^{D_\mathrm{fofe}}.
\end{equation}
By construction this is permutation-invariant and agnostic to the number of
visible entities.

\paragraph{Actor network (per role).}
The self-observation $o^{\text{self}}\in\mathbb R^{6}$ (ego kinematics + alive
flag) bypasses FOFE and is mapped by an MLP. The three entity channels are
encoded independently and concatenated with the self-embedding before a
fusion MLP and a factorised action head:
\begin{equation}
z = \big[\,g_\xi(o^{\text{self}})\,\big\Vert\,\mathrm{FOFE}_\text{ag}\,\big\Vert\,
        \mathrm{FOFE}_\text{tg}\,\big\Vert\,\mathrm{FOFE}_\text{rd}\,\big],\qquad
\bm\ell = W_{\text{head}}\,\mathrm{FusionMLP}(z).
\end{equation}

\paragraph{Critic network (per role).}
The critic receives the decomposed \emph{global} state: all agents
($d{=}7$: pos, speed, heading, $\omega$, role, alive), all targets ($d{=}3$),
all radars ($d{=}4$: pos, active flag, effective range), plus a scalar
time feature $\tau\in[0,1]$:
\begin{equation}
V_{\phi_\rho}(s) = W_v\;\mathrm{FusionMLP}\!\Big(
\big[\,\mathrm{FOFE}_\text{ag}(s)\,\big\Vert\,\mathrm{FOFE}_\text{tg}(s)\,\big\Vert\,
      \mathrm{FOFE}_\text{rd}(s)\,\big\Vert\,\tau\,\big]\Big).
\end{equation}

\paragraph{Adaptations from standard FOFE.}
\begin{enumerate}[leftmargin=1.25em]
\item \emph{Per-channel decomposition.} Wang's original encoder processes a
single heterogeneous entity stream. We split the observation into
semantically typed channels (agents, targets, radars) each with an
independent FOFE block, and treat the ego state as a separate fixed-size
MLP branch. This enforces a type-aware prior and makes channel-level
diagnostics (dominance, collapse rate) tractable.
\item \emph{Centralised FOFE critic with time feature.} The critic
consumes the full global state as three entity sets \emph{plus} an explicit
normalised-time scalar $\tau$, enabling the baseline to discriminate early
from late episode phases without extending the set with a dummy entity.
\item \emph{Masked max-pool with all-masked fallback.} When no entity in a
channel is visible ($\sum_i m_i=0$) we substitute a zero vector rather
than propagating $-\infty$ from the masked max, ensuring stable gradients
under aggressive occlusion.
\item \emph{Role-specific FOFE weights.} Striker and jammer each own a full
set of three FOFE blocks (actor and critic); no weight sharing across roles.
\end{enumerate}

\subsection{Reward Structure}
\label{subsec:rewards}
Let $d^{\max}_{ij}$ denote a component-specific cutoff distance and
$w^{\text{lin}}_{\bullet}$ the linear weight. All active shaping terms in our
configuration reduce to the linear sub-case of the more general
piecewise linear--exponential family (exponential weights set to zero):
\begin{equation}
\varphi(d;w,d^{\max})=w\cdot\max\!\left(0,\;1-\tfrac{d}{d^{\max}}\right).
\end{equation}

\paragraph{Sparse team rewards with team-spirit mixing.}
For a team event $e\in\{\text{kill target},\text{agent destroyed}\}$ with
scalar value $R_e$ and individual attribution indicator
$\mathbb{1}_{i=\text{credit}(e)}$, agent $i$'s received reward is
\begin{equation}
r_i^{e} = \eta \cdot \frac{R_e}{|\mathcal N^\text{alive}|}
       + (1-\eta)\cdot R_e\,\mathbb{1}_{i=\text{credit}(e)},
\qquad \eta=0.5,
\end{equation}
with $R_{\text{target}}=+3$, $R_{\text{destroy}}=-3$. A one-shot
\emph{terminal bonus} $R_{\text{term}}=+1$ is granted to every alive agent
at the transition where the mission is completed (all targets killed).

\paragraph{Per-step dense terms.}
For each alive agent $i$ at step $t$:
\begin{align}
r^{\text{step}}_{i,t} &= -0.01, \\
r^{\text{border}}_{i,t} &= -\varphi\big(d^{\text{edge}}_{i,t};\,0.05,\,0.05\big), \\
r^{\text{radar-av}}_{i,t} &= -\sum_k \varphi\big(d^{\text{rz}}_{i,k,t};\,0.05,\,0.1\big), \\
r^{\text{appr}}_{i,t} \;(\text{strikers only})
   &= \varphi\big(d^{\text{tgt}}_{i,t};\,0.1,\,1.0\big)-0.1, \\
r^{\text{form}}_{i,t} \;(\text{jammers only})
   &= 0.1\!\cdot\!\Big(\max\!\big(0,\,1-\tfrac{d^{\text{str}}_{i,t}}{0.5}\big)-1\Big), \\
r^{\text{sep}}_{i,t}  &= -\varphi\big(d^{\text{same-role}}_{i,t};\,0.05,\,0.1\big), \\
r^{\text{ctrl}}_{i,t} &= -0.01\,a_{i,t}^{2} - 0.01\,\alpha_{i,t}^{2},
\end{align}
where $d^{\text{edge}}$ is the distance to the nearest map edge,
$d^{\text{rz}}_{i,k}$ the distance to the boundary of the (fixed, jammed) zone
of radar $k$, $d^{\text{tgt}}_{i}$ the distance to the nearest alive target,
$d^{\text{str}}_{i}$ and $d^{\text{same-role}}_{i}$ the distances to the
nearest cross-role and same-role teammate respectively, and
$a,\alpha\in[-1,1]$ the discrete accel/angular-accel multipliers.
The striker-approach shaping is sign-flipped to act as a \emph{distance
penalty} (zero at $d{=}0$, monotonically decreasing with distance), which
preserves the same gradient magnitude as a conventional approach reward
while removing the positive bias that interfered with the sparse
$R_{\text{target}}$ signal.

\paragraph{Per-agent reward.}
Combining the above, each agent's per-step reward is
\begin{equation}
r_{i,t} = r^{e}_{i,t} + r^{\text{step}}_{i,t}
+ r^{\text{border}}_{i,t} + r^{\text{radar-av}}_{i,t}
+ r^{\text{appr}}_{i,t} + r^{\text{form}}_{i,t}
+ r^{\text{sep}}_{i,t} + r^{\text{ctrl}}_{i,t}
+ R_{\text{term}}\,\mathbb{1}_{\text{terminal}},
\end{equation}
evaluated only over alive agents and passed per-role into the advantage
computation of the corresponding MAPPO learner.

\paragraph{Adaptations from standard shaped-reward MAPPO.}
(i) The \emph{team-spirit} mixing~(12) interpolates between a fully
egoistic credit-assignment scheme ($\eta=0$) and a fully cooperative one
($\eta=1$); $\eta=0.5$ is used to balance individual responsibility with
joint objective pursuit. (ii) Dense shaping that conventionally takes the
form of positive approach bonuses is reformulated as a zero-at-contact
distance penalty~(16), which removes a positive reward baseline that
otherwise dominates the sparse engagement signal. (iii) The
same-role separation penalty~(18) and control-effort penalty~(19) act as
mild regularisers promoting spatial diversity and smooth trajectories
without biasing the sparse task reward.

\end{document}
